\scanpage{3}
$n\geqslant1$ fixé, et $n$-hyperpolyèdre combinatoire épinglé universel $ {}_n\Pi$, groupe d'autom.
\[
\unsure{\mathscr G_n|\mathfrak g_n|G_n|\mathfrak G_n}=\{\tau_0,\ldots,\tau_n\mid \tau_i\longleftrightarrow\tau_j\text{ si }0\leq i,j\leq n,\ i,j\text{ non cons.}\}
\]
$S$ schéma de base fixé

\margin{autodual} $C_{pr}$ cat. des réal. géom. régulières \inserted{non dég} projectives de $\Pi$ (autodual)

\margin{autodual} $C'_{pr}$ cat des systèmes $\deleted{(X,\varphi_G:\mathscr G_n\to\operatorname{Aut}_S X)}$ $(X,(\tau_i)_{0\leq i\leq n},\deleted{\unsure{b}}X_*)$ $X$ fibré projectif \inserted{rg $n+1$}, $\tau_i$ $(0\leq i\leq n)$ ps. refl. strictes \inserted{($\Rightarrow h_i,H_i$)}, \deleted{$s_0$ section} $X_*$ drapeau maximal, avec conditions
\[
\begin{aligned}
\text{a)}\ &\tau_i\longleftrightarrow\tau_j\qquad\text{si}\qquad 0\leq i,j\leq n\qquad i,j\text{ non consécutifs.}\\
\text{b)}\ &\tau_i X_j=X_j\quad\text{si }0\leq i\ne j\leq n\\
\text{c)}\ &\tau_i X_i\ne X_i\qquad 0\leq i\deleted{\unsure{,j}}\leq n
\end{aligned}
\]
\margin{\uline{NB} $C_{pr},C'_{pr}$ proviennent par \uline{rigidification} de catégories $\widetilde C_{pr}$ \deleted{\unsure{$C'_{pr}$}} de \uline{représentations} projectives de $\mathscr G_n$ \deleted{sans} plus}
\margin{\uline{NB} Au lieu de se donner $X_*$, il suffit de se donner \uline{un} $X_i$ avec $\tau_jX_i=X_i$ $(j\ne i)$ $\tau_iX_i\ne0$ satisfaisant à conditions ouv. -- p. ex $s_0=X_0$, tel que \deleted{la suite des \unsure{$s$}} \inserted{$(s_0,(h_i))$} engendre \unsure{lin$^{\mathrm{t}}$} $X$}

\margin{non autodual} $C_{aff}$ variante affine de $C_{pr}$

\margin{non autodual} $C'_{aff}$ variante affine de $C'_{pr}$ \inserted{[NB. il n'est plus besoin de se donner $\varphi_G$ \unsure{; idem|; itou} pour $C'_{aff}$].}

(Attention, $\tau_0$ n'est peut-être pas une ps. réflexion affine, mais $\widehat\tau_0$ est une ps. refl. projective (dont le \deleted{hyp} centre peut être exceptionnellement l'hyperplan à l'$\infty$ \ldots\ )
\margin{cas $\beta_0=\mu_0=0$ ! (avéré en car. $2$..$)$}

\margin{autodual} $C_{vect}$ variante vectorielle de $C_{pr}$ (où $X$ \uline{rigidisé} comme $P(E)$, et $s_0$ rigidisé par une section \uline{de} $E$ -- ce qui \deleted{rigidifie} est une donnée de \uline{rigidification} dans une donnée $\in$ Catégorie $\widetilde C_{vect}$, éliminant le groupe $G_m$ d'autom..
\margin{$\widetilde C_{vect},\widetilde C'_{vect}$ sans $G_m$-épinglage}

$C'_{vect}$ \unsure{variante} vectorielle de $C'_{pr}$ (on laisse tomber \unsure{qui est|qu'est} $=$ la \unsure{der. prem.|drapeau}, remplacée par $s_0=a_{-1}\ldots$)

$\deleted{\unsure{\,\sim\,}}C^{\infty}_{\unsure{pr*|pr\ne *}}$ cat des $(C,\varphi_G^C:\mathscr G_n\to\operatorname{Aut} C,b_0)$ où $C$ fibré projectif de rg $n$, $\varphi_G^C$ hom groupe, $h_0$ section de $C$, avec
\margin{non autodual, sauf si $\unsure{\rho|p}\ne\unsure{0|h_0}$ i.e. $c\notin C$ en tte fibre}
\margin{\uline{NB} La donnée de $s_0$ superflue si $\tau_0^C$ stricte i.e. $\mu_0,\beta_0$ non nuls simultanément}
a) les $\tau_i^C$ ps. réflex. $(0\leq i\leq n)$ strictes si $1\leq i\leq n$, associées à $\unsure{h_0|b_0}$ pour $\tau_0^C$ (qui peut être non stricte ($\Rightarrow \deleted{\unsure{\exists}}h_i,H_i$ pour $1\leq i\leq n$ -- mais pas de $h_0$)

% archive sheet number: 2

